\documentclass[11pt]{article}
\usepackage{amsmath, amssymb, amsthm}
\usepackage{geometry}
\geometry{margin=1in}
\usepackage{hyperref}

\newtheorem{theorem}{Theorem}
\newtheorem{proposition}[theorem]{Proposition}
\newtheorem{lemma}[theorem]{Lemma}
\newtheorem{corollary}[theorem]{Corollary}
\newtheorem{remark}[theorem]{Remark}

\title{The odd-$c$ constant leaf of the SHARED-CONTENT LEMMA:\\
$\mathrm{content}(G_4^{(c)}) = 4$ for odd $c \ge 7$}
\author{Clio}
\date{2026-07-21}

\begin{document}
\maketitle

\begin{abstract}
Fix an odd integer $c \ge 7$ and let $G_4^{(c)}(a,b)$ be the polynomial in
$\mathbb{Z}[a,b]$ obtained from the three-row engine
$M_j = \langle s_{(a,b,c)}, h_1^{2(m-j)}e_2^{\,j}\rangle$ at $j=4$ by peeling
off the $2(c-5)$ ``lower runs'' (definition below, matching the even-$c$
convention of \cite{clio2026capidentity}). We prove that
\[
  \mathrm{content}(G_4^{(c)}) \;=\; 4,
\]
where content is the minimum $2$-adic valuation of the binomial-basis
coefficients on the two odd-parity sheets $(\varepsilon_a, \varepsilon_b) \in
\{(0,1), (1,0)\}$. The proof is a finite case check on the $50$ coefficient
polynomials of $G_4^{(c)}$ in $\mathbb{Z}[c]$: every coefficient has $v_2 \ge 4$
uniformly in odd $c$, and the coefficient $P^{(1,0)}_{1,1}(c) = 24\,c\,(25c^3 -
114c^2 + 131c - 12)$ has $v_2 = 4$ exactly, for every odd $c$.
\end{abstract}

\section{Setup}

Fix a three-row partition $(a,b,c)$ with $a > b > c \ge 7$, $c$ odd, and
$a + b + c$ even (so that $|(a,b,c)| = 2m$ for an integer $m$). Since $c$ is
odd, $a + b$ is odd, so $a$ and $b$ have \emph{opposite} parities.

The three-row engine polynomial at index $j$ is
\begin{equation}\label{eq:Mj}
  M_j \;=\; \langle s_{(a,b,c)},\; h_1^{2(m-j)} e_2^{\,j}\rangle
        \;=\; \binom{2(m-j)}{b-j}\,\frac{(a-b+1)\,Q_c(a,b,j)}{c!\,(a+c+1-j)\prod_{i=1}^{c}(b+i-j)}.
\end{equation}
The heavy--tip decomposition of $Q_c$ is
\begin{equation}\label{eq:Qc}
  Q_c(a,b,j) \;=\; (a-(c-2))(b-(c-1))\,H_c(a,b,j) \;-\; (-1)^c\,(2c)!\,\binom{j}{2c}.
\end{equation}
At $j = 4$ and $c \ge 3$, we have $4 < 2c$ so $\binom{j}{2c} = 0$, and
$H_c(a,b,4)$ is the polynomial part of $Q_c(a,b,4)/((a-(c-2))(b-(c-1)))$.
By the degree count of \cite{clio2026capidentity}, $H_c(a,b,4)$ has bi-degree
$(c-1,\,c-1)$ in $(a,b)$. Dividing by the ``lower runs'' peels off $2(c-5)$
linear factors:
\begin{equation}\label{eq:G4}
  G_4^{(c)}(a,b) \;:=\; \frac{H_c(a,b,4)}{\displaystyle\prod_{s=3}^{c-3}(a+s)\cdot \prod_{s=2}^{c-4}(b+s)}.
\end{equation}

\begin{proposition}[Biquartic normal form; \cite{clio2026capidentity}]
\label{prop:biquartic}
For every integer $c \ge 6$, $G_4^{(c)}(a,b)$ is a polynomial in
$\mathbb{Z}[a,b]$ of bi-degree $(4,4)$. Its coefficients (as functions of $c$)
are polynomials in $\mathbb{Z}[c]$ of degree at most $4$.
\end{proposition}

The proof of Proposition~\ref{prop:biquartic} in \cite{clio2026capidentity}
does not use the parity of $c$: the bi-degree $(4,4)$ statement follows from
the degree count $(c-1) - (c-5) = 4$ (which holds for all $c \ge 6$), and
integrality of the lower-runs quotient was verified there by symbolic
interpolation for $c \in \{6, 8, \ldots, 26\}$ and confirmed cross-computed
against direct $M_j$ evaluation. For the odd-$c$ range relevant here
($c \in \{7, 9, \ldots, 31\}$), integrality was verified explicitly during the
present work.

\subsection{The content on the parity sheet}\label{ssec:content}

For odd $c$, the parity sheets on which $s_{(a,b,c)}$ is a genuine
Littlewood--Richardson coefficient (i.e., $a + b + c$ even) are indexed by
$(\varepsilon_a, \varepsilon_b) \in \{(0,1), (1,0)\}$: substitute
\[
  a = 2u + \varepsilon_a, \qquad b = 2v + \varepsilon_b, \qquad (u,v) \in \mathbb{Z}_{\ge 0}^2.
\]
Under this substitution, $G_4^{(c)}(2u+\varepsilon_a, 2v+\varepsilon_b)$ is a
polynomial in $(u,v)$ of bi-degree $(4,4)$ with integer coefficients, admitting
a unique binomial-basis expansion
\[
  P(u,v) \;=\; \sum_{i=0}^{4}\sum_{k=0}^{4} \alpha_{i,k}\,\binom{u}{i}\binom{v}{k},
  \qquad \alpha_{i,k} = (\Delta_u^{\,i} \Delta_v^{\,k} P)(0,0).
\]
The $2$-adic content on the parity sheet is
\[
  \mathrm{cont}_2\bigl(G_4^{(c)} \text{ on sheet }(\varepsilon_a, \varepsilon_b)\bigr)
    \;:=\; \min_{i,k}\, v_2(\alpha_{i,k})
    \;=\; \min_{(u,v)\in\mathbb{Z}^2} v_2\bigl(G_4^{(c)}(2u+\varepsilon_a, 2v+\varepsilon_b)\bigr).
\]
The second equality is the standard Mahler / integer-Newton statement. We
define
\[
  \mathrm{content}(G_4^{(c)}) \;:=\; \min\Bigl\{\mathrm{cont}_2 \bigl(G_4^{(c)} \text{ on sheet } (\varepsilon_a, \varepsilon_b)\bigr) : (\varepsilon_a, \varepsilon_b) \in \{(0,1),(1,0)\}\Bigr\}.
\]

\section{Theorem statement}

\begin{theorem}[Odd-$c$ constant leaf]\label{thm:main}
For every odd integer $c \ge 7$,
\[
  \mathrm{content}(G_4^{(c)}) \;=\; 4.
\]
\end{theorem}

The upper bound (Section~\ref{sec:upper}) exhibits a single universal witness
that has $v_2 = 4$ exactly for every odd $c$. The lower bound
(Section~\ref{sec:lower}) is a finite check that every one of the $50$
coefficient polynomials has $v_2 \ge 4$ uniformly.

\section{The coefficient polynomials}\label{sec:coeffs}

Let $c \ge 7$ be odd. Fix a parity sheet
$(\varepsilon_a, \varepsilon_b) \in \{(0,1), (1,0)\}$ and set
\[
  P^{(\varepsilon_a, \varepsilon_b)}_{i,k}(c) \;:=\; (\Delta_u^{\,i} \Delta_v^{\,k})\bigl[G_4^{(c)}(2u+\varepsilon_a, 2v+\varepsilon_b)\bigr]_{(u,v)=(0,0)}.
\]
By Proposition~\ref{prop:biquartic}, each is a polynomial in $\mathbb{Z}[c]$
of degree at most $4$, obtained by interpolation from the values at, e.g.,
$c \in \{7, 9, 11, 13, 15\}$ (and cross-verified at $c \in \{17, 19, \ldots, 31\}$,
$0$ mismatches).

\subsection{Sheet $(\varepsilon_a, \varepsilon_b) = (0,1)$}

\begin{small}
\begin{align*}
(0,0)&: 48\,c(c{-}1) & (2,3)&: 2304(2c^3 + 18c^2 + 32c + 25) \\
(0,1)&: 24\,c(c{-}1)(5c^2 - 37c + 64) & (2,4)&: 9216(2c^2 + 6c + 5) \\
(0,2)&: 24\,c(c{-}1)(25c^2 - 121c + 150) & (3,0)&: 96\,c(c{-}2)(c{-}1)(9c - 7) \\
(0,3)&: 96\,c(c{-}2)(c{-}1)(9c - 7) & (3,1)&: 384\,c(c^3 + 16c^2 + 2c + 11) \\
(0,4)&: 384\,c(c{-}2)(c{-}1)(c{+}1) & (3,2)&: 2304(2c^3 + 18c^2 + 32c + 25) \\
(1,0)&: 24\,c(c{-}1)(5c^2 - 37c + 64) & (3,3)&: 9216(3c^2 + 18c + 25) \\
(1,1)&: 24\,c(25c^3 - 106c^2 + 117c - 44) & (3,4)&: 36864(2c + 5) \\
(1,2)&: 96\,c(9c^3 + 16c^2 - 36c + 33) & (4,0)&: 384\,c(c{-}2)(c{-}1)(c{+}1) \\
(1,3)&: 384\,c(c^3 + 16c^2 + 2c + 11) & (4,1)&: 1536\,c(c{+}1)(2c + 1) \\
(1,4)&: 1536\,c(c{+}1)(2c + 1) & (4,2)&: 9216(2c^2 + 6c + 5) \\
(2,0)&: 24\,c(c{-}1)(25c^2 - 121c + 150) & (4,3)&: 36864(2c + 5) \\
(2,1)&: 96\,c(9c^3 + 16c^2 - 36c + 33) & (4,4)&: 147456 \\
(2,2)&: 192(2c^4 + 44c^3 + 90c^2 + 100c + 75) & &
\end{align*}
\end{small}

\subsection{Sheet $(\varepsilon_a, \varepsilon_b) = (1,0)$}

\begin{small}
\begin{align*}
(0,0)&: 0 & (2,3)&: 2304(2c^3 + 20c^2 + 44c + 39) \\
(0,1)&: 24\,c(c{-}1)(5c^2 - 36c + 58) & (2,4)&: 9216(2c^2 + 10c + 13) \\
(0,2)&: 24\,c(c{-}1)(25c^2 - 93c + 74) & (3,0)&: 288\,c(c{-}3)(c{-}2)(3c - 5) \\
(0,3)&: 288\,c(c{-}1)(3c^2 - c + 2) & (3,1)&: 384\,c(c^3 + 14c^2 - 31c + 19) \\
(0,4)&: 384\,c(c{-}1)(c{+}1)(c{+}2) & (3,2)&: 2304(2c^3 + 16c^2 + 12c + 7) \\
(1,0)&: 120\,c(c{-}4)(c{-}3)(c{-}2) & (3,3)&: 27648(c^2 + 6c + 7) \\
(1,1)&: 24\,c(25c^3 - 114c^2 + 131c - 12) & (3,4)&: 36864(2c + 7) \\
(1,2)&: 96(9c^4 + 22c^3 - 28c^2 + 21c + 6) & (4,0)&: 384\,c(c{-}3)(c{-}2)(c{-}1) \\
(1,3)&: 384(c^4 + 18c^3 + 29c^2 + 39c + 18) & (4,1)&: 1536\,c(c{-}1)(2c - 1) \\
(1,4)&: 1536(c{+}1)(c{+}2)(2c + 3) & (4,2)&: 9216(2c^2 + 2c + 1) \\
(2,0)&: 120\,c(c{-}3)(c{-}2)(5c - 13) & (4,3)&: 36864(2c + 3) \\
(2,1)&: 96\,c(9c^3 + 6c^2 - 68c + 73) & (4,4)&: 147456 \\
(2,2)&: 192(2c^4 + 44c^3 + 78c^2 + 52c + 39) & &
\end{align*}
\end{small}

\section{Basic $2$-adic facts for odd $c$}\label{sec:facts}

For $c$ odd and any integer polynomial $P(c) \in \mathbb{Z}[c]$, the reduction
$P(c) \bmod 2^k$ depends only on $c \bmod 2^k$ (since $P(c_1) - P(c_2)$ is
divisible by $(c_1 - c_2)$). This lets us determine $v_2(P(c))$ for all odd $c$
by evaluating $P$ modulo a fixed power of $2$, using representatives $c = 1, 3,
5, 7, \ldots$ up to the required residue class.

We use the following elementary facts throughout (all for $c$ odd,
$c \ge 7$):

\begin{lemma}\label{lem:parity}
For every odd integer $c$:
\begin{enumerate}
\item[(i)]  $v_2(c) = v_2(c-2) = v_2(c-4) = v_2(c+2) = 0$.
\item[(ii)] $v_2(c-1) \ge 1$, $v_2(c-3) \ge 1$, $v_2(c-5) \ge 1$, $v_2(c+1) \ge 1$.
\item[(iii)] If $P(c) \in \mathbb{Z}[c]$ has odd constant term and every other
             coefficient even, then $v_2(P(c)) = 0$ for all $c$.
\item[(iv)] $c(c-1) \equiv 0 \pmod 2$ and, more sharply, $v_2(c(c-1)) \ge 1$.
\end{enumerate}
\end{lemma}

\begin{proof}
(i)--(ii) are direct: $c$ odd forces $c-2$ and $c-4$ odd (they differ from
$c$ by an even number) and $c-1, c-3, c-5, c+1$ even. (iii) is the standard
parity observation. (iv) follows from (i)--(ii).
\end{proof}

\section{Upper bound}\label{sec:upper}

\begin{lemma}[Universal witness]\label{lem:upper}
For every odd integer $c \ge 7$,
\[
  v_2\bigl(P^{(1,0)}_{1,1}(c)\bigr) \;=\; 4.
\]
Consequently, $\mathrm{content}(G_4^{(c)}) \le 4$.
\end{lemma}

\begin{proof}
From the coefficient table in Section~\ref{sec:coeffs},
$P^{(1,0)}_{1,1}(c) = 24\,c\,(25c^3 - 114c^2 + 131c - 12)$. Write
\[
  R(c) \;:=\; 25c^3 - 114c^2 + 131c - 12.
\]
Then $v_2(P^{(1,0)}_{1,1}(c)) = v_2(24) + v_2(c) + v_2(R(c)) = 3 + 0 + v_2(R(c))$
(using $c$ odd), so it suffices to show $v_2(R(c)) = 1$ for every odd $c$.

Reduce $R$ modulo $4$. For $c$ odd, $c^2 \equiv 1 \pmod 8$ (write $c = 2t+1$;
$(2t+1)^2 = 4t(t+1) + 1$ and $t(t+1)$ is even), hence $c^2 \equiv 1 \pmod 4$
and $c^3 = c \cdot c^2 \equiv c \pmod 4$. Compute each term mod $4$:
\begin{align*}
25 c^3 &\equiv c \pmod 4, & \text{since } 25 \equiv 1 \pmod 4, \\
114 c^2 &\equiv 114 \cdot 1 \equiv 2 \pmod 4, \\
131 c &\equiv 3c \pmod 4, & \text{since } 131 \equiv 3 \pmod 4, \\
12 &\equiv 0 \pmod 4.
\end{align*}
Therefore
\[
  R(c) \;\equiv\; c - 2 + 3c - 0 \;=\; 4c - 2 \;\equiv\; 2 \pmod 4,
\]
so $v_2(R(c)) = 1$ for every odd $c$. Thus $v_2(P^{(1,0)}_{1,1}(c)) = 3 + 1 = 4$.
\end{proof}

\section{Lower bound}\label{sec:lower}

\begin{proposition}\label{prop:lower}
For every odd integer $c \ge 7$ and every $(\varepsilon_a, \varepsilon_b) \in
\{(0,1),(1,0)\}$ and every $(i,k) \in \{0,1,2,3,4\}^2$,
\[
  v_2\bigl(P^{(\varepsilon_a,\varepsilon_b)}_{i,k}(c)\bigr) \;\ge\; 4.
\]
Consequently, $\mathrm{content}(G_4^{(c)}) \ge 4$.
\end{proposition}

We prove this by inspection of each of the $50$ coefficients. To avoid
repetition, we group them by the mechanism that provides the lower bound.

\subsection{High-prefactor coefficients}\label{ssec:high-prefactor}

For each coefficient $P^{(\cdot)}_{i,k}(c) = 2^\alpha \cdot Q(c)$ with
$Q(c) \in \mathbb{Z}[c]$, if $\alpha \ge 4$ then $v_2(P) \ge 4$ trivially. This
handles the following on \emph{both} sheets (reading prefactors from
Section~\ref{sec:coeffs}):
\begin{center}
$\alpha \ge 5$: (0,3), (0,4), (1,2), (1,3), (1,4), (2,1), (2,2), (2,3), (2,4),$\!$
(3,0), (3,1), (3,2), (3,3), (3,4), (4,0), (4,1), (4,2), (4,3), (4,4).
\end{center}
These are $19$ index pairs on each of the $2$ sheets, i.e., $38$ coefficients.
The remaining $12$ coefficients (with prefactor $2^3 \cdot m$ or $2^4 \cdot 3$,
or the identically-zero coefficient $P^{(1,0)}_{0,0}$) require a targeted
argument, given below.

\subsection{Coefficients with prefactor $24 = 2^3 \cdot 3$}\label{ssec:24}

There are $8$ coefficients with prefactor $24$; they come in two shapes.

\paragraph{Shape A: $24\, c\, (c-1) \cdot Q(c)$ where $Q(c)$ is a quadratic with parity-preserving structure.}
On sheet $(0,1)$: $(0,1), (1,0), (0,2), (2,0)$.
On sheet $(1,0)$: $(0,1), (0,2)$.

For $c$ odd, $v_2(24) + v_2(c) = 3 + 0 = 3$ and $v_2(c-1) \ge 1$
(Lemma~\ref{lem:parity}(ii)), so it suffices to show $v_2(Q(c)) \ge 0$
(automatic since $Q(c) \in \mathbb{Z}$).

Sharper: for these coefficients, $v_2(Q(c))$ is easily determined:
\begin{itemize}
\item Sheet $(0,1)$, $(0,1) = (1,0)$: $Q = 5c^2 - 37c + 64$. For $c$ odd:
$5c^2$ odd, $37c$ odd, $64$ even $\Rightarrow$ $Q \equiv \text{odd} + \text{odd} + \text{even} \equiv \text{even} \pmod 2$. So $v_2(Q) \ge 1$, giving total $v_2 \ge 5$.
\item Sheet $(0,1)$, $(0,2) = (2,0)$: $Q = 25c^2 - 121c + 150$. For $c$ odd:
odd $+$ odd $+$ even $\equiv$ even. $v_2(Q) \ge 1$; total $\ge 5$.
\item Sheet $(1,0)$, $(0,1)$: $Q = 5c^2 - 36c + 58$. For $c$ odd:
$5c^2$ odd, $36c$ even, $58$ even $\Rightarrow$ $Q$ odd. $v_2(Q) = 0$; total
$= 3 + v_2(c-1) \ge 3 + 1 = 4$.
\item Sheet $(1,0)$, $(0,2)$: $Q = 25c^2 - 93c + 74$. For $c$ odd:
odd $+$ odd $+$ even $\equiv$ even. $v_2(Q) \ge 1$; total $\ge 5$.
\end{itemize}
All $6$ coefficients of shape A satisfy $v_2 \ge 4$.

\paragraph{Shape B: $24\, c \cdot R(c)$ where $R(c)$ is a cubic.}
On sheet $(0,1)$: $(1,1)$.
On sheet $(1,0)$: $(1,1)$.

For $c$ odd, $v_2(24) + v_2(c) = 3$, so we need $v_2(R(c)) \ge 1$.

\begin{itemize}
\item Sheet $(0,1)$, $(1,1)$: $R(c) = 25c^3 - 106c^2 + 117c - 44$.
Compute mod $4$ (using $c^2 \equiv 1$, $c^3 \equiv c$ mod $4$ for $c$ odd):
\[
25c^3 - 106c^2 + 117c - 44 \;\equiv\; c - 2 + c - 0 \;=\; 2c - 2 \pmod 4.
\]
For $c$ odd, $2c - 2 = 2(c-1)$ with $c-1$ even, so $2c - 2 \equiv 0 \pmod 4$.
Hence $v_2(R) \ge 2$; total $v_2 \ge 3 + 2 = 5$.
\item Sheet $(1,0)$, $(1,1)$: $R(c) = 25c^3 - 114c^2 + 131c - 12$; handled in
Lemma~\ref{lem:upper}, where we showed $R \equiv 2 \pmod 4$ hence $v_2(R) = 1$.
Total $v_2 = 4$. (This is the tight witness.)
\end{itemize}

\subsection{Coefficients with prefactor $48 = 2^4 \cdot 3$}\label{ssec:48}

Only one: $P^{(0,1)}_{0,0}(c) = 48\, c\, (c-1)$. Since $v_2(48) = 4$, total
$v_2 \ge 4$. $\checkmark$

\subsection{Coefficients with prefactor $96 = 2^5 \cdot 3$}\label{ssec:96}

Already handled by Section~\ref{ssec:high-prefactor}.

\subsection{Coefficients with prefactor $120 = 2^3 \cdot 15$}\label{ssec:120}

There are two, both on sheet $(1,0)$: $(1,0)$ and $(2,0)$.

\paragraph{Sheet $(1,0)$, $(1,0)$: $P = 120\, c\, (c-4)(c-3)(c-2)$.}
For $c$ odd, $v_2(120) + v_2(c) + v_2(c-4) + v_2(c-3) + v_2(c-2) =
3 + 0 + 0 + v_2(c-3) + 0 \ge 3 + 1 = 4$. $\checkmark$

\paragraph{Sheet $(1,0)$, $(2,0)$: $P = 120\, c\, (c-3)(c-2)(5c-13)$.}
For $c$ odd, $5c - 13 = \text{odd} - \text{odd} = \text{even}$, so $v_2(5c-13) \ge 1$;
also $v_2(c-3) \ge 1$. Total $v_2 \ge 3 + 0 + 1 + 0 + 1 = 5$. $\checkmark$

\subsection{Coefficients with prefactor $288 = 2^5 \cdot 9$}\label{ssec:288}

Handled by Section~\ref{ssec:high-prefactor} (prefactor $2^5$).

\subsection{Verification summary}

Every one of the $50$ coefficient polynomials
$P^{(\varepsilon_a, \varepsilon_b)}_{i,k}(c)$ has $v_2 \ge 4$ for all odd
$c \ge 7$:

\begin{itemize}
\item $38$ coefficients are handled by the high-prefactor argument
      (Section~\ref{ssec:high-prefactor}): prefactor divisible by $2^5$
      ($19$ on each sheet).
\item $1$ coefficient has prefactor $2^4 \cdot 3 = 48$ (Section~\ref{ssec:48}).
\item $2$ coefficients have prefactor $120 = 2^3 \cdot 15$
      (Section~\ref{ssec:120}), gaining the missing factor of $2$ from $c-3$
      (odd $c$) or from $5c-13$ (odd $c$).
\item $6$ coefficients of Shape A (Section~\ref{ssec:24}), prefactor $24$, gain
      $v_2 \ge 1$ from $c-1$ and (in five of six cases) additional $v_2 \ge 1$
      from the inner polynomial being even at odd $c$.
\item $2$ coefficients of Shape B (Section~\ref{ssec:24}), prefactor $24$,
      require the mod-$4$ analysis: sheet-$(1,0)$ $(1,1)$ gains exactly $v_2 = 1$
      (giving total $4$), sheet-$(0,1)$ $(1,1)$ gains $v_2 \ge 2$ (giving total
      $\ge 5$).
\item $1$ identically-zero coefficient $P^{(1,0)}_{0,0} = 0$ (contributes
      $\infty$ and does not lower the content).
\end{itemize}
Total: $38 + 1 + 2 + 6 + 2 + 1 = 50$ index pairs, accounting for all
coefficients on both sheets.

This proves Proposition~\ref{prop:lower}.

\section{Proof of Theorem~\ref{thm:main}}

Combining Lemma~\ref{lem:upper} (upper bound $\mathrm{content} \le 4$) and
Proposition~\ref{prop:lower} (lower bound $\mathrm{content} \ge 4$):
\[
  \mathrm{content}(G_4^{(c)}) \;=\; 4 \qquad \text{for every odd integer } c \ge 7. \qed
\]

\section{Computational verification}\label{sec:verify}

\begin{enumerate}
\item \textbf{Extended empirical anchor.} $\mathrm{content}(G_4^{(c)}) = 4$ was
      verified numerically for all odd $c \in \{7, 9, 11, \ldots, 99\}$ ($47$
      values) by two independent engines: (i) Lyra's Pieri +
      Jacobi--Trudi route (uid 469, verified odd $c \in \{7, \ldots, 31\}$
      with double-anchor $G_4^{(7)}(2, 1) = 51408 = 2^4 \cdot 3213$), and
      (ii) the direct $M_4$ evaluator
      \texttt{compute/2026-07-21-odd-c-content-extension/direct\_M4.py} (odd $c \in \{33, \ldots, 99\}$),
      cross-checked $0/157$ mismatch on three-row shapes with $|\lambda| \in [10, 22]$.
\item \textbf{Coefficient polynomials.} The $50$ polynomials in
      Section~\ref{sec:coeffs} were interpolated from the $5$ values
      $c \in \{7, 9, 11, 13, 15\}$ and cross-verified at
      $c \in \{17, 19, 21, 23, 25, 27, 29, 31\}$: $0/400$ mismatch across the
      $50 \times 8 = 400$ interpolant-vs-actual comparisons.
\item \textbf{Coefficient-level $v_2$.} Every coefficient's minimum $v_2$ over
      odd $c \in \{7, 9, \ldots, 99\}$ agrees with the case-analysis bound in
      Section~\ref{sec:lower}. The tight witness $P^{(1,0)}_{1,1}(c)$ evaluates
      to $v_2 = 4$ at every odd $c$ in the range, as predicted by the
      Lemma~\ref{lem:upper} mod-$4$ computation.
\end{enumerate}

\section{Consequences}

\begin{corollary}[Physical $\Delta(4)$ for odd $c$]
Using the physical identity $\Delta(4) = 2\,v_2(G_4^{(c)}) - 10$
(memory node \texttt{2026-07-05-generalc-gen4-crux-proved}), we get
\[
  \min \Delta(4) \;=\; 2 \cdot 4 - 10 \;=\; -2 \qquad \text{for every odd } c \ge 7,
\]
i.e., the generator $j = 4$ contributes a negative shift of $2$ units below
the tip-run scale on odd sheets. Contrast with the even-$c$ leaf
(\cite{clio2026capidentity}): $\min \Delta(4) \in \{0, 2\}$ depending on
$c \bmod 4$.
\end{corollary}

\begin{remark}[Sheet asymmetry and the $-1$ offset]
The odd sheets are asymmetric: $(1,0)$ has content $4$ while $(0,1)$ has
content $5$. The $K$-witness of the even-$c$ leaf,
$K(c) = 24\,c(c-1)(c-4)(c-5)$, has $v_2 \ge 3 + 0 + 1 + 0 + 1 = 5$ for odd
$c$; naively this would predict content $= 5$. The observed content $4$
therefore signals a genuinely new witness on sheet $(1,0)$, and
Lemma~\ref{lem:upper} identifies it: the cubic factor $25c^3 - 114c^2 +
131c - 12$ inside the $(1,1)$ coefficient. Concretely, this cubic reduces to
$2 \pmod 4$ for every odd $c$ (a single-bit residue). The ``$-1$ offset''
observed in the numerics is precisely the single-bit cancellation
$v_2(K)_{\text{even}} - v_2(P^{(1,0)}_{1,1}) \ge 5 - 4 = 1$: the tight sheet
$(1,0)$ carries a coefficient that beats the $K$-witness by exactly one bit,
uniformly in odd $c$.
\end{remark}

\begin{remark}[Toward the SHARED-CONTENT LEMMA]
Theorem~\ref{thm:main} together with \cite{clio2026capidentity} completes
Instance $3$ of the SHARED-CONTENT LEMMA in the two-tower form:
\[
  \mathrm{content}(G_4^{(c)}) \;=\; \begin{cases}
    \min(v_2(K(c)),\, 6) & \text{if } c \text{ even, } c \ge 6, \\
    4 & \text{if } c \text{ odd, } c \ge 7.
  \end{cases}
\]
Instance $1$ (Clio $\clubsuit$ resonance, Lyra's lane) and Instance $2$
(uncapped two-tower, Lyra's $\gamma(c)$) remain open at the time of writing.
\end{remark}

\section*{Acknowledgments}

Lyra independently rebuilt $G_4^{(c)}$ via Pieri and reported
$\mathrm{content}(G_4^{(c)}) = 4$ constantly on odd $c \in \{7, \ldots, 31\}$
(uid 469, 2026-07-21), double-anchored at $c = 7$ via $51408 = 2^4 \cdot
3213$; that observation triggered the present PROVE cycle. The background
probe extending the range to odd $c \in \{33, \ldots, 99\}$ (147 partitions
cross-checked, $0$ mismatch) sharpened the target from a suspected ``$c = 7$
hole'' into a clean constant leaf. Rick's abacus-blueprint suggestion (uid
454) previously motivated the finite-substitution style of proof used here.

\begin{thebibliography}{9}

\bibitem{clio2026capidentity}
Clio, \emph{Lyra's cap identity for $G_4^{(c)}$: $\mathrm{content}(G_4^{(c)}) = \min(v_2(K(c)), 6)$ for even $c \ge 6$},
proof file \texttt{proofs/2026-07-20-lyra-cap-identity.tex}, 2026.

\bibitem{clio2026generalc}
Clio, \emph{General-$c$ master valuation for the three-row heavy engine},
memory node \texttt{2026-06-17-generalc-master-and-c5}, 2026.

\end{thebibliography}

\end{document}
